
%%% ====================================================================
%%% ====================================================================
%%%
%%%  APPENDIX: FORMALIZATION AND VERIFICATION
%%%
%%% ====================================================================
%%% ====================================================================
\appendix

\section{Lean Formalization Audit}\label{app:lean}

This section documents the formal verification of the algebraic chain
from $\golden^2=\golden+1$ to the spectral identification with~$\zeta$.

\subsection{Audit summary}

\begin{center}
\begin{tabular}{lrrrl}
\toprule
Source File & Lines & Decls. & Sorry & Status \\
\midrule
\lean{odd_zeta_pcf\_v1.lean} & 658 & 67 & 0 & Verified \\
\lean{verify_odd_zeta_pcf\_unified.py} & 905 & 220 & --- & All pass \\
\bottomrule
\end{tabular}
\end{center}

The Lean file has a single axiom: \lean{phi\_sq} ($\golden^2 = \golden + 1$).
All other results are derived from this axiom, Mathlib, and decidable computation.

\subsection{Deductive chain}

The formalization is organized into four parts:

\begin{center}
\begin{tabular}{cll}
\toprule
Part & Content & Key identifiers \\
\midrule
1 & PCF isomorphism chain & \lean{pentagon\_identity}, \lean{G\_hom}, \\
  & $\golden \to G \to \mu_3 \to \omega^k/2$ & \lean{spectral\_uniqueness}, \lean{eigenvalue\_half} \\
2 & Euler product of $\zeta$ & \lean{primes_land_in\_Z20star}, \\
  & primes $\to$ $\chi_5$ $\to$ $f_p$ & \lean{fib_split\_criterion}, \lean{euler_type_from\_G} \\
3 & Transport & \lean{spectral\_identification} (\texttt{rfl}), \\
  & PCF spectral $=$ $\zeta$ spectral & \lean{zeta_positive_spectral\_half} \\
4 & $\mathbb{F}_1$, Weil, physics & \lean{F1_descent\_data}, \lean{weil\_analogy}, \\
  & descent, invariants & \lean{physical\_invariants}, \lean{discriminant\_five} \\
\bottomrule
\end{tabular}
\end{center}

\subsection{Complete lean audit}\label{ssec:lean-audit-table}

\begin{table}[htb]
\centering
\caption{Lean audit: correspondence between paper results and
formalised identifiers. All dependencies are fully verified.
The single axiom is $\golden^2=\golden+1$.}
\label{tab:lean-audit-app}
\footnotesize
\begin{tabular}{lll}
\toprule
Paper result & Lean identifier & Method \\
\midrule
\multicolumn{3}{l}{\textbf{Part 1: PCF isomorphism chain}} \\
$\golden^2=\golden+1$ & \lean{phi\_sq} & axiom \\
$4(\golden/2)^2-2(\golden/2)-1=0$ & \lean{phi_half\_quadratic} & \texttt{nlinarith} \\
$\cos(5\theta)$ expansion & \lean{cos_five\_mul} & \texttt{nlinarith} \\
$\cos(\pi/5)>0$ & \lean{cos_pi_five\_pos} & monotonicity \\
$4\cos^2(\pi/5)-2\cos(\pi/5)-1=0$ & \lean{cos_pi_five\_quadratic} & unique root \\
$\golden = 2\cos(\pi/5)$ & \lean{pentagon\_identity} & \texttt{linarith} \\
$\vert\mathbb{Z}_{20}^\times\vert=8$ & \lean{Z20star\_card} & \texttt{decide} \\
Mult.\ closed & \lean{mul\_closed} & \texttt{decide} \\
Add.\ destroyed & \lean{add\_destroyed} & \texttt{decide} \\
$G$ homomorphism & \lean{G\_hom} & \texttt{decide} \\
$G$ surjective & \lean{G\_surj} & \lean{fin\_cases} \\
$\vert\ker G\vert=2$ & \lean{G\_kernel} & \texttt{decide} \\
Exponent $2$ & \lean{exponent\_two} & \texttt{decide} \\
$\mu_2/\mu_3 = 2 = \vert\ker G\vert$ & \lean{contraction_is\_kernel} & \lean{norm\_num} \\
$(\sigma,\mu)=(3/2,1/2)$ unique & \lean{spectral\_uniqueness} & \texttt{nlinarith} \\
Diagonal blocked & \lean{diagonal\_blocked} & \texttt{nlinarith} \\
$\mathrm{Re}(\omega_{\mathrm{pcf}})=-1/2$ & \lean{omega\_properties} & \lean{exp\_mul\_I} \\
$\mathrm{Re}>0 \implies \mathrm{Re}=1/2$ & \lean{eigenvalue\_half} & exhaustion \\
\midrule
\multicolumn{3}{l}{\textbf{Part 2: Isomorphism with Euler product}} \\
$\forall p>5$: $p\bmod 20\in\mathbb{Z}_{20}^\times$ & \lean{primes\_land\_in\_Z20star} & \texttt{omega+decide} \\
$\chi_5$ from $G$ & \lean{chi5\_from\_G} & \texttt{rfl} \\
$\chi_5$ classification & \lean{chi5\_values} & \texttt{decide} \\
$\chi_5(p)\equiv F_p\pmod{p}$ & \lean{fib_split\_criterion} & \lean{native\_decide} \\
$f_p$ type from $G$ & \lean{euler_type_from\_G} & \texttt{rfl} \\
\midrule
\multicolumn{3}{l}{\textbf{Part 3: Transport}} \\
$\zeta_{\mathrm{spectral}} = \zeta_{\mathrm{PCF}}$ & \lean{spectral\_identification} & \texttt{rfl} \\
Transport principle & \lean{transport\_spectral} & rewrite \\
$\zeta > 0 \implies \mathrm{Re}(s)=1/2$ & \lean{zeta\_positive\_spectral\_half} & transport \\
\midrule
\multicolumn{3}{l}{\textbf{Part 4: $\mathbb{F}_1$-descent, Weil, physics}} \\
$z^2 = zy + y^2$ & \lean{self\_measurement} & \texttt{nlinarith} \\
$\psi_p\circ\psi_q = \psi_{pq}$ & \lean{frobenius\_composition} & \lean{pow\_mul} \\
$\mathbb{F}_1$-descent data & \lean{F1\_descent\_data} & conjunction \\
Weil analogy & \lean{weil\_analogy} & conjunction \\
$\mu_3^2 = 1/4$ & \lean{mu3\_squared} & \lean{norm\_num} \\
$1-\mu_3^2 = 3/4$ & \lean{holographic\_ratio} & \lean{norm\_num} \\
$c = 3\ell/(2G_N) = 3$ & \lean{brown\_henneaux\_c} & \lean{norm\_num} \\
$1/(4G_N) = \mu_3$ & \lean{BH\_factor} & \lean{norm\_num} \\
$-1/i = i$ (S-duality) & \lean{S\_duality\_fixed} & \lean{Complex.I\_mul\_I} \\
$\ell=1/\ell \implies \ell=1$ & \lean{ads\_radius\_one} & \texttt{nlinarith} \\
$\bar\golden = 1-\golden$ & \lean{phi\_bar\_eq} & \lean{field\_simp} \\
$\mathrm{Tr}(\golden) = 1$ & \lean{trace\_one} & \texttt{ring} \\
$N(\golden) = -1$ & \lean{norm\_neg\_one} & \lean{field\_simp} \\
$\mathrm{disc}(\mathbb{Q}(\sqrt{5})) = 5$ & \lean{discriminant\_five} & \texttt{nlinarith} \\
\bottomrule
\end{tabular}
\end{table}

\subsection{Numerical verification}

The Python script \lean{verify_odd_zeta_pcf\_unified.py} performs
220 independent checks at 50-digit precision covering:
Frobenius lifts (\cref{sec:background}),
prime decomposition (\cref{sec:primes}),
Euler product convergence,
factorisation (\cref{sec:factorisation}),
the fundamental formula (\cref{sec:fundamental}),
the even Bernoulli formula (\cref{sec:even}),
odd-value determination (\cref{sec:pentagon}),
the categorical framework (\cref{sec:categorical-framework}),
the Weil analogy (\cref{ssec:weil-analogy}),
the polylogarithmic circuit (\cref{ssec:polylog-circuit}),
and explicit $\mathbb{F}_1$-descent verification (Borger criteria).

\subsection{Scope of formalization}

The categorical structures of \cref{sec:categorical-framework}---the
Frobenius functor (functoriality verified via
$\psi_p\circ\psi_q=\psi_{pq}$), the co-determination theorem
(reducing to \lean{G\_hom}, \lean{G\_surj}, \lean{G\_kernel}),
and the spectral uniqueness (\lean{spectral\_uniqueness})---are
verified through the Lean formalization.

The Euler product colimit (\cref{prop:euler-colimit}),
the Dedekind natural isomorphism (\cref{thm:dedekind-natural}),
and the class-number formula are not currently formalisable:
Mathlib lacks Dirichlet series, the spectral mapping theorem
for Banach algebras, and the identity theorem for analytic continuation.
These are verified numerically (220 checks, all pass).

\section{Axiomatic basis: core equations}\label{app:axiomatic-basis}

The framework shares foundational equations ($\golden^2=\golden+1$,
$\mu_3=1/2$, $\tau_{\mathrm{PCF}}=i$, $M_{\mathrm{PCF}}$,
$\Lambda_{\mathrm{PCF}}$, and the Galois functor~$G$) with the
Primitive Structure~\cite{V11}.  The equations below are specific to
the arithmetic formulation developed in this paper, each established
in the indicated section and verified in \cref{sec:verification}.

\subsection*{Arithmetic of \texorpdfstring{$\mathbb{Q}(\sqrt{5})$}{Q(sqrt 5)}}
\begin{align}
  \bar\golden &= -1/\golden = 1 - \golden \label{eq:ax-phi-bar} \\
  \mathrm{Tr}(\golden) &= \golden + \bar\golden = 1 \label{eq:ax-trace} \\
  N(\golden) &= \golden \cdot \bar\golden = -1 \label{eq:ax-norm} \\
  \mathrm{disc}(\mathbb{Q}(\sqrt{5})) &= (\golden - \bar\golden)^2 = 5 \label{eq:ax-disc}
\end{align}

\subsection*{Frobenius and splitting}
\begin{align}
  \psi_p(\golden) &= \golden^p = F_p\,\golden + F_{p-1} \label{eq:ax-frobenius} \\
  \chi_5(p) &\equiv F_p \pmod{p} \label{eq:ax-fib-split} \\
  \chi_5 &= \pi_2 \circ G \label{eq:ax-chi5-G}
\end{align}

\subsection*{Mersenne bridge}
\begin{align}
  \lambda_{\log} &= \frac{\log 2}{\log\golden} \label{eq:ax-lambda-log} \\
  \golden^{\lambda_{\log}} &= 2 \label{eq:ax-mersenne} \\
  T &= 2\pi\log\golden \label{eq:ax-T}
\end{align}

\subsection*{Zeta and L-values}
\begin{align}
  \zeta_{\mathbb{Q}(\sqrt{5})}(s) &= \zeta(s)\cdot L(s,\chi_5) \label{eq:ax-dedekind} \\
  L(1,\chi_5) &= \frac{2\log\golden}{\sqrt{5}} = \frac{T}{\pi\sqrt{5}} \label{eq:ax-L1} \\
  \zeta(2k+1) &= \frac{\zeta_{\mathbb{Q}(\sqrt{5})}(2k+1)}{L(2k+1,\chi\_5)} \label{eq:ax-odd-zeta}
\end{align}

\subsection*{Physical invariants from \texorpdfstring{$\mu\_3=1/2$}{mu\_3=1/2}}
\begin{align}
  G_N &= \mu_3 = 1/2 \label{eq:ax-GN} \\
  \ell &= 1 \quad\text{(AdS radius from S-duality)} \label{eq:ax-ell} \\
  c &= 3\ell/(2G_N) = 3 \quad\text{(Brown--Henneaux)} \label{eq:ax-c} \\
  R &= \log\golden \quad\text{(regulator)} \label{eq:ax-R} \\
  S_{\mathrm{BH}}/k_B &= \lambda_{\log}\cdot R = \log 2 \label{eq:ax-SBH}
\end{align}

\section{Compilation and verification}\label{app:compilation}

The project requires Lean~4 with Mathlib. To verify the formal proofs:
\begin{enumerate}
  \item Navigate to the \texttt{lean/} directory and run \texttt{lake build}.
  \item Alternatively, from the root folder, run \texttt{pnpm run verify}
        to execute the dual verification suite (Lean and Python).
\end{enumerate}

All \texttt{decide} and \lean{fin\_cases} tactics terminate in
${}<30$~seconds on standard hardware.
Source files and numerical verification suite are available at:
\begin{center}
  \url{\pcfrepo} \\
  \url{https://omega-pcf.com/odd-zeta}
\end{center}
